% answers_short.tex —— 2015 年·全国Ⅱ卷·理科 · 简短答案速查
% 由 tools/make_short_answers.py 自 answers.tex 的 \gksolhead/\gkanswer 提取，
% 并与 2015/2015_数学真题_全国2理_答案要点.md 逐题比对后生成，勿手改（改源 MD 后重跑）。

\begin{gknotice}
\textbf{使用说明}：本节只给最终答案便于快速核对，完整推导见解析版。答案为编写者独立推导，\textbf{非官方标准答案}。
\end{gknotice}

\gkpart{选择题}
\begin{center}\begin{tabular}{ccccccccccccc}
\toprule
题号 & 1 & 2 & 3 & 4 & 5 & 6 & 7 & 8 & 9 & 10 & 11 & 12 \\
\midrule
答案 & A & B & D & B & C & D & C & B & C & B & D & A \\
\bottomrule
\end{tabular}\end{center}

\gkpart{填空题}
\begin{itemize}[leftmargin=2.6em]
  \item[\textbf{13.}] $\frac12$
  \item[\textbf{14.}] $8$
  \item[\textbf{15.}] $\frac{x^2}{4}-y^2=1$
  \item[\textbf{16.}] $a=8$
\end{itemize}

\gkpart{解答题}
\begin{itemize}[leftmargin=2.6em]
  \item[\textbf{17.}] (1) $\frac12$；(2) $BD=\sqrt2, AC=1$
  \item[\textbf{18.}] (1) A 地均值较高、分散较小；(2) $P(C)=\frac{12}{25}$
  \item[\textbf{19.}] (1) $E,F,Q,P$ 为截面正方形顶点，$AP=DQ=10$；(2) $\sin\beta=\frac4{3\sqrt5}$
  \item[\textbf{20.}] (1) 斜率乘积定值 $-9$；(2) 能，$k=4\pm\sqrt7$
  \item[\textbf{21.}] (1) 负半轴递减、正半轴递增；(2) $m\in[-1,1]$
  \item[\textbf{22.}] (1) $EF\parallel BC$；(2) 面积 $\frac{16\sqrt3}{3}$
  \item[\textbf{23.}] (1) $(0,0),(\frac{\sqrt3}{2},\frac32)$；(2) 最大值 $4$，$\alpha=\frac{5\pi}{6}$
  \item[\textbf{24.}] (1) 命题成立；(2) 两条件互为充要条件
\end{itemize}
